\scanpage{18}
\mspage{5}
récurrence, il faut prouver a) i.e. \deleted{\unsure{l'iso}} que
\deleted{$\underline{\lim} X_i\to\unsure{\lim}$}
\[
\begin{aligned}
u_\alpha\overset{\mathrm{déf}}{=}\varinjlim_{i<\alpha}\deleted{\unsure{X}}u_i:\quad&
 \underset{\inserted{\substack{\mathrm{déf}\ \Vert\\\varinjlim_{i<\alpha}X_i}}}{X_\alpha}
 \unsure{-|\to}
 \underset{\inserted{\substack{\Vert\\\varinjlim_{i<\alpha}Y_i}}}{Y_\alpha}
\end{aligned}
\]
est dans $C$. \deleted{Et} Quant \inserted{à c, il est déjà} \deleted{\unsure{si on suppose, en} plus de $u_.\in C_{I_\alpha}$, que $u_.\in W_{I^\circ}$, donc (par hyp. de réc.) que pour tt $i_0<\alpha$, dans $(\alpha)$ on a
\[
 \widetilde{u}_{i_0}\in TC\overset{\mathrm{déf}}{=}C\cap W
\]
(et non seulement $\widetilde{u}_{i_0}\in C$), alors on a $u_\alpha\in$}
contenu dans l'hyp. de récurrence.

Donc il reste finalement à prouver
\[
 u_\alpha\in C
\]
quand on suppose que la prop. est prouvée pour les ordinaux $\beta<\alpha$.

On va construire un troisième système \unsure{ordinal|ordonné} sur $I_\alpha$, en \deleted{\unsure{posant}} introduisant pour tt $0\leq i<\alpha$, \inserted{i.e. $i\in I_\alpha$,} $Z_i$ par le carré cocartésien
\diagram{\begin{tikzcd}[ampersand replacement=\&]
 X_i\arrow[r,-]\arrow[d]\arrow[dr,phantom,"\text{cocart}" description] \& X_\alpha\arrow[d]\\
 Y_i\arrow[r] \& Z_i
\end{tikzcd}}
les flèches de transition entre les $Z_i$ sont évidentes, et on aura \deleted{\unsure{c'est}} $\varinjlim_i Z_{\deleted{\unsure{\alpha}}i}\simeq Y_\alpha$. Ainsi $u_\alpha:X_\alpha\inserted{(=Z_{-1})}-Y_\alpha$ \deleted{\unsure{se}} est devenu un composé transfini de morphismes \deleted{\unsure{$Z_i\to Z_{i+1}$} \unsure{$(i\leq\alpha)$} \illegible{1}}
% archive sheet number: 17

